% !TITLE 茅屋为秋风所破歌
% !AUTHOR 杜甫
% !DYNASTY 唐
% !ORDER 22

\begin{poem}
八月秋高风怒号，卷我屋上三重茅\textsuperscript{1}。
茅飞渡江洒江郊，高者挂罥长林梢\textsuperscript{2}，下者飘转沉塘坳\textsuperscript{3}。

南村群童欺我老无力，忍能对面为盗贼\textsuperscript{4}。
公然抱茅入竹去，唇焦口燥呼不得，归来倚杖自叹息。

俄顷风定云墨色\textsuperscript{5}，秋天漠漠向昏黑\textsuperscript{6}。
布衾多年冷似铁\textsuperscript{7}，娇儿恶卧踏里裂\textsuperscript{8}。
床头屋漏无干处，雨脚如麻未断绝\textsuperscript{9}。
自经丧乱少睡眠\textsuperscript{10}，长夜沾湿何由彻\textsuperscript{11}！

安得广厦千万间\textsuperscript{12}，大庇天下寒士俱欢颜\textsuperscript{13}！
风雨不动安如山。
呜呼！何时眼前突兀见此屋\textsuperscript{14}，吾庐独破受冻死亦足\textsuperscript{15}！
\end{poem}

\begin{annotations}
\notetext{1}{三重茅：好几层茅草。三，虚数，多的意思。狂风把屋顶厚厚的茅草卷走了。}
\notetext{2}{挂罥长林梢：挂（缠绕）在高高的树梢上。罥（juàn），缠绕、悬挂。}
\notetext{3}{塘坳：池塘的低洼处。坳（ào），低凹的地方。}
\notetext{4}{忍能对面为盗贼：竟然忍心当面做贼。孩子们抱着茅草跑进竹林，杜甫喊得口干舌燥也叫不住。}
\notetext{5}{俄顷风定云墨色：一会儿风停了，乌云墨黑。俄顷，片刻。}
\notetext{6}{漠漠向昏黑：天色阴沉沉地暗了下来。漠漠，阴沉迷蒙的样子。}
\notetext{7}{布衾多年冷似铁：布被子用了多年，硬邦邦的，像铁一样冷。衾（qīn），被子。}
\notetext{8}{娇儿恶卧踏里裂：孩子睡相不好，把被里蹬破了。恶卧，睡相不安。}
\notetext{9}{雨脚如麻未断绝：雨点像麻线一样密密麻麻，下个不停。雨脚，雨丝。}
\notetext{10}{丧乱：战乱，指安史之乱。自从经历战乱就很少能睡好觉。}
\notetext{11}{何由彻：怎么熬到天亮。“彻”即彻晓、到天亮。}
\notetext{12}{广厦：宽敞的大屋。安得，怎样才能得到。}
\notetext{13}{大庇天下寒士：庇护天下所有贫寒的读书人。庇（bì），遮蔽、保护；寒士，贫寒的士人。}
\notetext{14}{突兀见此屋：忽然在眼前出现这样的房子。突兀（wù），突然高耸而出。}
\notetext{15}{吾庐独破受冻死亦足：即便只有我的茅屋破漏、受冻至死，也心满意足了。这是杜甫推己及人的最高境界。}
\end{annotations}

\begin{appreciation}
这首诗写于上元二年（761年），杜甫漂泊到成都，在浣花溪边盖了一座茅屋。那年秋天，一场大风把屋顶的茅草卷走了，紧接着又下起大雨，全家人在漏雨的屋子里挨了一整夜。但杜甫没有止于自怜，而是从个人的苦难中跳了出来，喊出了那句震撼千古的呼喊。

八月秋高风怒号，卷我屋上三重茅——八月深秋，狂风怒号，卷走了我屋顶上三层茅草。三重茅不是说真的有三层，而是形容茅草之多。茅飞渡江洒江郊，高者挂罥长林梢，下者飘转沉塘坳——茅草飞过江，洒落在江郊；高的挂在树梢上，低的飘转着沉到池塘里。挂罥（juàn）就是缠绕悬挂。这几句写得很生动，像慢镜头一样展示了茅草被风卷走的全过程。

南村群童欺我老无力，忍能对面为盗贼——南村的孩子们欺负我年老无力，竟然忍心当面做贼，抱着我的茅草跑进了竹林。这里杜甫用了盗贼一词，似乎有些过激。但紧接着他就写道：唇焦口燥呼不得，归来倚杖自叹息——喊得唇焦口燥也没有用，只好回来拄着拐杖叹气。他不是真的在恨那些孩子，他是在恨自己的无能为力。

布衾多年冷似铁，娇儿恶卧踏里裂——被子用了多年，冷得像铁；孩子睡相不好，把被里都蹬破了。这是一个穷苦人家的真实写照。床头屋漏无干处，雨脚如麻未断绝——床上没有一块干的地方，雨点像麻线一样密密麻麻，下个不停。自经丧乱少睡眠，长夜沾湿何由彻——自从经历战乱，就很少能睡安稳；这一夜又湿又冷，怎么才能熬到天亮？

然后，诗忽然转到了一个完全不同的境界。安得广厦千万间，大庇天下寒士俱欢颜——怎样才能有千万间宽大的房屋，庇护天下所有贫寒的读书人，让他们都露出笑脸？这不是一个反问，是一个愿望。杜甫从自己一家人的苦难，想到了天下所有像他一样的人。他的茅屋破了，但他想的不是自己，是天下寒士。

呜呼！何时眼前突兀见此屋，吾庐独破受冻死亦足——什么时候眼前能突然出现这样的房屋，即使我的茅屋独破、受冻而死，也心满意足了。这是一种几乎宗教式的献身精神。杜甫不是圣人，他也在乎自己的冷暖；但他在乎天下人胜过在乎自己。这种推己及人的胸怀，是中国文化中最珍贵的品质。
\end{appreciation}
